% Proposed isolated A.94. No frozen cycle30 source is changed.
\statement{A.94}{Proved}{The full affine-character boundary of the marked pentagon}
Use the standard pure generators
\[
       (a,b,c,d,e,f)=(A_{12},A_{13},A_{23},A_{14},A_{24},A_{34}).
\tag{A.94.1}
\]
For every field k, every homomorphism
\(P_4\to\operatorname{Aff}_1(k)=k\rtimes k^\times\)
with nonabelian image factors through one of the following
five marked free-two quotients, expressed by the images of
the original generators:
\[
\begin{array}{c|l}
\mathrm{(i)}&a=b=d=1,\quad cef=1,\\
\mathrm{(ii)}&a=c=e=1,\quad bdf=1,\\
\mathrm{(iii)}&b=c=f=1,\quad ade=1,\\
\mathrm{(iv)}&d=e=f=1,\quad abc=1,\\
\mathrm{(v)}&a=f,\quad c=d,\quad b=(df)^{-1},\quad e=(fd)^{-1}.
\end{array}
\tag{A.94.2}
\]
Consequently, if \(u\in\widehat F(X,Y)'\) satisfies the
exact swap duality \(u(Y,X)=u(X,Y)^{-1}\), its marked pentagon
\[
       u(a,ce)u(bc,f)=u(c,f)u(ab,ef)u(a,c)
\tag{A.94.3}
\]
holds after every homomorphism from \(\widehat P_4\) to a
finite affine group over a field. More generally it holds after
every homomorphism to any finite group with an elementary
abelian two-kernel and odd abelian quotient.

Thus the entire first elementary-two layer over every odd
abelian head adds no pentagon condition beyond exact swap
duality. This is a theorem about these actual quotient maps,
not about the pentagon in \(\widehat P_4\), the realization
of canonical parameters, or the vanishing of the full cusp word.

\begin{proof}
\emph{The original marked source.}\par
The Fadell--Neuwirth decomposition is
\(P_4=F(d,e,f)\rtimes P_3\). In the inverse-base conjugation
convention its exact pure actions are
\[
\begin{array}{c|ccc}
 &d&e&f\\ \hline
 \operatorname{Ad}(a^{-1})&\operatorname{Ad}(de)d&\operatorname{Ad}(d)e&f\\
 \operatorname{Ad}(b^{-1})&\operatorname{Ad}(df)d&\operatorname{Ad}([d,f])e&\operatorname{Ad}(d)f\\
 \operatorname{Ad}(c^{-1})&d&\operatorname{Ad}(ef)e&\operatorname{Ad}(e)f.
\end{array}
\tag{A.94.4}
\]
Here \([d,f]=dfd^{-1}f^{-1}\). These are precisely the
four-case Artin formulas of \citet[Example~1.1,
pp.~465--466]{CECohen2010}, with its preceding right-action
convention, specialized to the standard generators in
(A.94.1). Only this actual action formula is used from that
source. In the base,
\[
 Z=abc\text{ is central},\qquad P_3=F(a,c)\times\langle Z\rangle,
 \qquad b=a^{-1}Zc^{-1}.
\]
Composing the three inverse actions, or using the full twists
\(\Delta_3^2=abc\), \(\Delta_4^2=abcdef\), gives
\[
          \operatorname{Ad}(Z^{-1})|_{F(d,e,f)}=operatorname{Ad}(def).
\tag{A.94.5}
\]
The complete a and c rows also follow directly from centrality
of the actual three-strand twists ade and cef, together with
\([a,f]=[c,d]=1\). Thus the interlaced b-row is not an
additional unproved marking assumption in the calculation below.
Equation (A.94.3) is literally
\citet[p.~547, Eq.~(4)]{Furusho2010} under the standard
pure-generator identification immediately following that
display. No prounipotent-to-profinite implication from that
paper is used.

Write affine multiplication as
\[
                 (v,t)(w,s)=(v+tw,ts).
\]
The centralizer of every nonidentity affine element is abelian.
A nonzero translation has the translation subgroup as its
centralizer; an element of slope different from one has a
unique fixed point and its centralizer is the multiplicative
group about that point. In particular the centralizer of a
nonabelian affine subgroup is trivial.

\emph{Reduction to a nonabelian fiber.}\par
Write D,E,F for the fiber images. If their subgroup is nontrivial
abelian, each action in (A.94.4) fixes every image. The base
therefore centralizes that subgroup, and the entire image lies
in the abelian centralizer of one of its nonidentity elements.
This contradicts the hypothesis. If the fiber is trivial,
the image factors through \(P_3\); its central Z must be
trivial because the image is nonabelian. This gives (iv).

Henceforth the fiber is nonabelian. Put P=DEF. By (A.94.5),
the images of Z and \(P^{-1}\) have the same conjugation
action on the fiber. The fiber centralizer is trivial, so
\[
                         Z=P^{-1}.
\tag{A.94.6}
\]
In particular both a and c commute with P.

\emph{The exact two rank-one matrices.}\par
Write
\[
\begin{gathered}
 D=(v_1,t_1),\quad E=(v_2,t_2),\quad F=(v_3,t_3),
 \quad\rho=t_1t_2t_3,\\
 v=(v_1,v_2,v_3)^t,\quad
 p=(1-t_1,1-t_2,1-t_3)^t,\\
 L(v)=v_1+t_1v_2+t_1t_2v_3,\qquad L(p)=1-\rho.
\end{gathered}
\tag{A.94.7}
\]
The translation of P is L(v). The a and c rows of (A.94.4)
act on the translation coordinates by
\[
\begin{gathered}
 A=I+u_A\ell_A,\quad u_A=(t_1,-1,0)^t,
       \quad\ell_A=(t_2-1,1-t_1,0),\\
 C=I+u_C\ell_C,\quad u_C=(0,t_2,-1)^t,
       \quad\ell_C=(0,t_3-1,1-t_2).
\end{gathered}
\tag{A.94.8}
\]
These follow directly from
\(\operatorname{Ad}(w,q)(v,t)=((1-t)w+qv,t)\).
Both rows kill p, and both columns lie in \(\ker L\).
If the image of a is \((w_a,a_0)\), the inverse-base
convention gives the exact equation
\[
                    Av=a_0^{-1}v-a_0^{-1}w_ap,
\tag{A.94.9}
\]
and similarly for c. In the basis \(u_A,u_C\) of \(\ker L\)
the matrices are
\[
 A=\begin{pmatrix}t_1t_2&t_2(1-t_1)\\0&1\end{pmatrix},
 \qquad
 C=\begin{pmatrix}1&0\\1-t_3&t_2t_3\end{pmatrix}.
\tag{A.94.10}
\]
The columns are independent because \(t_1,t_2,t_3\) are
nonzero. The determinant of their two fixed-vector equations is
\[
 (t_1t_2-1)(t_2t_3-1)-t_2(1-t_1)(1-t_3)
                 =(t_2-1)(\rho-1).
\tag{A.94.11}
\]
No root of unity or bounded residue is selected in these identities.

\emph{The case \(\rho\ne1\).}\par
Conjugate the entire image so that \(P=(0,\rho)\).
Its centralizer is the diagonal torus, so \(w_a=w_c=0\).
We have \(L(v)=0\), hence \(v=\xi u_A+\eta u_C\),
and v is nonzero, since otherwise the fiber would be abelian.
Equations (A.94.9)--(A.94.10) say that v is a common eigenvector.
If \(\eta=0\), then
\(Cu_A=u_A+(1-t_3)u_C\) forces \(t_3=1\); also \(v_3=0\),
so \(F=1\). If \(\xi=0\), the symmetric A equation gives
\(t_1=1\) and \(v_1=0\), so \(D=1\).

If \(\xi\eta\ne0\), the second coordinate of the A equation
and the first coordinate of the C equation force both
eigenvalues to be one. Equation (A.94.11) gives \(t_2=1\).
The remaining equations give \(\xi=\eta\), since at least
one of \(t_1,t_3\) differs from one. Thus
\(v_2=-\xi+t_2\eta=0\), so \(E=1\).

\emph{The case \(\rho=1\).}\par
First suppose P is a nonzero translation. Its centralizer
is the translation subgroup, so \(a_0=c_0=1\), and
\[
                 u_A\ell_A(v)\in kp,
                 \qquad u_C\ell_C(v)\in kp.
\tag{A.94.12}
\]
Not all \(t_i\) equal one, since that would make the fiber
a translation group. If none equals one, neither column is
proportional to p. Both row values must vanish. The two rows
are independent, as their first and last nonzero entries
show, so their common kernel is kp. Then \(L(v)=0\),
contrary to the nonzero translation of P.

Thus at least one \(t_i=1\). If \(t_3=1\), then
\(t_1t_2=1\) with both factors different from one.
The column \(u_C\) is not proportional to p, so
\(\ell_C(v)=(1-t_2)v_3=0\), giving \(F=1\).
If \(t_1=1\), the A row gives \(D=1\).
If \(t_2=1\), then \(t_1t_3=1\) with both different
from one; \(u_A\notin kp\) and
\(\ell_A(v)=(1-t_1)v_2=0\), giving \(E=1\).
Two slopes equal to one would force all three equal to one,
already excluded.

Finally suppose \(P=1\). Then \(E=(FD)^{-1}\).
The complete a-inverse action on the fiber is conjugation
by \(F^{-1}\), and the complete c-inverse action is
conjugation by \(D^{-1}\). Triviality of the fiber
centralizer gives \(a=F,c=D\). Equation (A.94.6) gives
\(Z=1\), and hence \(b=(DF)^{-1}\). This is (v).

\emph{Exhaustion and literal free-two quotients.}\par
In the cases above where \(F=1\), the b and c actions
are identity on the nonabelian fiber. Its trivial centralizer
gives \(b=c=1\); the a action gives \(a=(DE)^{-1}\).
This is (iii). If \(D=1\), similarly \(a=b=1\) and
\(c=(EF)^{-1}\), giving (i). If \(E=1\), one gets
\(a=c=1\) and \(b=(DF)^{-1}\), giving (ii).
All equalities are literal, not only conjugation identities.
Each of the five substitutions satisfies the complete
(A.94.4) and the P3 relations on two freely chosen generators.
Thus they are the specified actual free-two quotients, and
the proof exhausts every nonabelian affine image over every field.

\emph{The entire pentagon evaluation.}\par
An abelian image makes every u-value trivial. In each of
(i)--(iv), substitution into (A.94.3) leaves the same single
u-value on both sides; all other argument pairs generate
cyclic groups. The surviving values, respectively, are
\[
                    u(c,f),\quad u(b,f),\quad
                    u(a,e),\quad u(a,c).
\]
For (v), the two left argument pairs are \((f,f^{-1})\)
and \((f^{-1},f)\), and the middle right pair is
\((d^{-1},d^{-1})\). Those three values are trivial,
leaving
\[
                            1=u(d,f)u(f,d),
\]
the given full duality. Finite affine targets admit all
the profinite-word evaluations, so this proves the first
pentagon assertion without a finite-image census.

\emph{Every character of an odd abelian head.}\par
Let Q be finite abelian of odd order and V a finite
\(\mathbb F_2[Q]\)-module. Over a finite extension field
containing all roots of unity of exponent Q, each Q-operator
has square-free minimal polynomial, and the operators commute.
Successive eigenspace decompositions simultaneously diagonalize
V after scalar extension. Every scalar coordinate of every
character space gives a homomorphism
\[
                         V\rtimes Q\longrightarrow
                                      \operatorname{Aff}_1(k).
\]
These coordinates separate V. Since Q is abelian and u is
a commutator word, the discrepancy of (A.94.3) has trivial
Q-head. Every separating affine coordinate kills it by the
proved classification, so the discrepancy is trivial in
\(V\rtimes Q\).

A finite extension of Q by V also splits. Indeed the finite
homogeneous-cochain averaging proved internally in B.143,
applied to the trivial subgroup, gives \(|Q|c=\partial b\)
for every two-cocycle c with the actual Q-action. Multiplication
by the odd integer \(|Q|\) is invertible on V, so
\(c=\partial(|Q|^{-1}b)\). The section corrected by that
cochain is a homomorphism. Thus the preceding conclusion
applies to every such finite extension, independently of a
chosen splitting.

This retains all characters of all odd abelian heads at once.
It proves a precise limit of this first elementary-two pentagon
comparison. Further nonlinear kernel information, a different
head with retained action, or additional actual canonical
compatibility is required to go beyond it. No assertion that
arbitrary duality data are canonical, no full coefficient
equality, and no nonlinear \(\Omega_\alpha=1\) conclusion
is contained in the theorem.
\end{proof}
